% Generated draft: requires mathematical and source review.
\documentclass[11pt]{article}
\usepackage[margin=25mm]{geometry}
\usepackage{fontspec}
\setmainfont{Latin Modern Roman}
\usepackage{amsmath,amssymb,mathtools,unicode-math}
\setmathfont{Latin Modern Math}
\usepackage{enumitem,needspace,xurl,hyperref}
\setlength{\parindent}{0pt}
\setlength{\parskip}{5pt}
\setlength{\emergencystretch}{3em}
\newcommand{\N}{\mathbb{N}}
\newcommand{\Z}{\mathbb{Z}}
\newcommand{\Q}{\mathbb{Q}}
\newcommand{\R}{\mathbb{R}}
\newcommand{\C}{\mathbb{C}}
\newcommand{\F}{\mathbb{F}}
\newcommand{\A}{\mathbb{A}}
\newcommand{\PP}{\mathbb{P}}
\newcommand{\E}{\mathbb{E}}
\newcommand{\eps}{\varepsilon}
\newcommand{\dd}{\,\mathrm{d}}
\begin{document}
\section*{UnsolvedMath ID: 701179}
\textbf{Source problem code:} catalogue entry
\subsection{Definitions and mathematical statement}
For a field \(F\), an ordinary algebraic representation of a group \(G\) is a homomorphism \(\rho:G\to\operatorname{GL}_F(V)\), where \(V\) is a nonzero \(F\)-vector space. It is faithful if \(\ker\rho=\{1\}\), and irreducible if its only \(G\)-invariant \(F\)-linear subspaces are \(0\) and \(V\). For a subgroup \(H\leq G\) and an \(F[H]\)-module \(W\), ordinary algebraic induction is
\[
\operatorname{Ind}_H^G W=F[G]\otimes_{F[H]}W,
\]
with \(G\) acting by left multiplication on the first factor.

The supplied question, proposed by A. V. Tushev, is: Is every faithful irreducible representation of a linear group \(G\) of finite rank over a finitely generated field of characteristic zero induced from an irreducible representation of a finitely generated dense subgroup of \(G\)? Here a finitely generated field of characteristic zero is a field \(k=\Q(a_1,\ldots,a_m)\) for some finite tuple, and being linear over \(k\) means admitting an embedding \(G\hookrightarrow\operatorname{GL}_n(k)\) for some positive integer \(n\).

Under the interpretation of representations and induction given above, the requested conclusion has the form
\[
\exists\,H\leq G\;\exists\,W\neq0:\quad
H\text{ is finitely generated and dense in }G,\qquad
W\text{ is an irreducible }F[H]\text{-module},\qquad
V\cong F[G]\otimes_{F[H]}W
\]
as \(F[G]\)-modules, for every representation satisfying the source hypotheses. The field \(F\) here is a parameter for explaining algebraic induction, not an additional source hypothesis. The record does not specify the meaning of finite rank or dense, the coefficient field of the representations, or whether representations must be finite dimensional. It also does not unambiguously identify whether the finitely generated field qualifies the linear realization of \(G\), the representation under consideration, or both. These qualifications must be resolved before this schema becomes a fully specified mathematical statement.
\subsection{Short English statement}
Tushev asks whether every faithful irreducible representation of a linear group satisfying the stated finite-rank and characteristic-zero field conditions comes by induction from an irreducible representation of a finitely generated dense subgroup. The meanings of rank and density, and the precise representation-field and dimension conventions, remain unspecified in the supplied record.
\subsection{Sources}
\begin{itemize}
\item Supplied local UnsolvedMath registry record: permanent ID \(701179\), source code \texttt{catalogue entry}, title \emph{Kourovka Notebook Problem 15.91}; upstream revision \texttt{d371ee917ac9647fbb06cd6fc529deec6f3dda82}.
\item Supplied attribution: A. V. Tushev, proposed in 2002. The record identifies the source as \emph{The Kourovka Notebook (Unsolved Problems in Group Theory)}, 15th issue (2002), Problem 15.91, and the 21st issue edition (2026), \texttt{arXiv:1401.0300v48}.
\item The supplied collection description attributes editorship to E. I. Khukhro and V. D. Mazurov. The registry labels this problem open; that label and the bibliographic details are supplied source data and have not been independently verified.
\end{itemize}
\end{document}
